\documentclass[11pt]{article}

\usepackage{amsmath, amssymb, amsthm}
\usepackage{geometry}
\usepackage{enumitem}
\usepackage{setspace}

\geometry{margin=1in}
\onehalfspacing

\title{\textbf{Lecture 6 Scribe: Random Variables and Distributions}}
\author{}
\date{}

\begin{document}
\maketitle

\section*{(1) Previous Lecture Recap: Random Variables}

\subsection*{Definition of a Random Variable}
A \textbf{Random Variable (RV)} is a function that maps outcomes of a random experiment to real numbers.

\[
X: \Omega \rightarrow \mathbb{R}
\]

\textbf{Reasoning:}  
The outcome of a random experiment is often non-numerical (e.g., heads/tails). To perform mathematical analysis such as computing averages or probabilities, we assign numbers to outcomes. This assignment is exactly what a random variable does.

\subsection*{Types of Random Variables}
\begin{itemize}
    \item \textbf{Discrete Random Variable}: Takes countable values.
    \item \textbf{Continuous Random Variable}: Takes values over an interval.
\end{itemize}

In this lecture, we focus primarily on \textbf{discrete random variables}.

\section*{(2) Independent Events}

\subsection*{Definition}
Two events $A$ and $B$ are said to be \textbf{independent} if
\[
P(A \cap B) = P(A)P(B)
\]

\textbf{Reasoning:}  
Independence means the occurrence of one event does not affect the probability of the other. If knowing $A$ does not change the likelihood of $B$, then the joint probability should factor into the product of their individual probabilities.

\subsection*{Example}
If a coin is tossed twice, let:
\[
A = \{\text{first toss is Head}\}, \quad B = \{\text{second toss is Head}\}
\]
Then,
\[
P(A) = \frac{1}{2}, \quad P(B) = \frac{1}{2}, \quad P(A \cap B) = \frac{1}{4}
\]
Since $P(A \cap B) = P(A)P(B)$, the events are independent.

\section*{(3) Types of Discrete Random Variables}

\subsection*{Bernoulli Random Variable}

A random variable $X$ is a \textbf{Bernoulli RV} if:
\[
X =
\begin{cases}
1 & \text{with probability } p \\
0 & \text{with probability } 1-p
\end{cases}
\]

\textbf{Reasoning:}  
This RV models a single trial with only two outcomes: success or failure. Assigning values $1$ and $0$ simplifies probability and expectation calculations.

\subsection*{Binomial Random Variable}

A random variable $X$ is \textbf{Binomial} with parameters $(n, p)$ if it counts the number of successes in $n$ independent Bernoulli trials.

\[
P(X = k) = \binom{n}{k} p^k (1-p)^{n-k}, \quad k = 0,1,\dots,n
\]

\textbf{Reasoning:}  
\begin{itemize}
    \item $\binom{n}{k}$ counts the number of ways to choose $k$ successes.
    \item $p^k$ is the probability of $k$ successes.
    \item $(1-p)^{n-k}$ is the probability of $n-k$ failures.
\end{itemize}

Multiplying gives the probability of exactly $k$ successes.

\subsection*{Geometric Random Variable}

A random variable $X$ is \textbf{Geometric} if it represents the number of trials until the first success.

\[
P(X = k) = (1-p)^{k-1} p, \quad k = 1,2,\dots
\]

\textbf{Reasoning:}  
The first $k-1$ trials must fail, and the $k$th trial must succeed. Independence allows multiplication of probabilities.

\subsection*{Poisson Random Variable}

A random variable $X$ follows a \textbf{Poisson distribution} with parameter $\lambda > 0$ if:
\[
P(X = k) = \frac{\lambda^k e^{-\lambda}}{k!}, \quad k = 0,1,2,\dots
\]

\textbf{Reasoning:}  
This distribution models the number of occurrences of rare events over a fixed interval, where events occur independently at a constant average rate.

\section*{(4) Expectation of Random Variables}

\subsection*{Definition}
For a discrete RV $X$:
\[
\mathbb{E}[X] = \sum_x x P(X = x)
\]

\textbf{Reasoning:}  
Expectation is a weighted average of all possible values, where weights are their probabilities.

\subsection*{Example}
If $X$ is Bernoulli$(p)$:
\[
\mathbb{E}[X] = 1 \cdot p + 0 \cdot (1-p) = p
\]

\subsection*{Expectation of a Function of an RV}
For a function $g(X)$:
\[
\mathbb{E}[g(X)] = \sum_x g(x) P(X = x)
\]

\textbf{Reasoning:}  
Instead of transforming the distribution, we directly weight function values by original probabilities.

\subsection*{Linearity of Expectation}
For constants $a, b$:
\[
\mathbb{E}[aX + b] = a\mathbb{E}[X] + b
\]

\textbf{Reasoning:}  
Expectation distributes over addition and scalar multiplication regardless of independence.

\subsection*{Moments and Central Moments}

\textbf{Variance}:
\[
\text{Var}(X) = \mathbb{E}[(X - \mathbb{E}[X])^2]
\]

\textbf{Reasoning:}  
Variance measures spread by averaging squared deviations from the mean.

Higher-order central moments define:
\begin{itemize}
    \item \textbf{Skewness}: asymmetry
    \item \textbf{Kurtosis}: tail heaviness
\end{itemize}

\section*{(5) Cumulative Distribution Function (CDF)}

\subsection*{Definition}
The \textbf{CDF} of a random variable $X$ is:
\[
F_X(x) = P(X \le x)
\]

\textbf{Reasoning:}  
The CDF accumulates probability up to a given point, making it useful for probability queries.

\subsection*{Properties}
\begin{itemize}
    \item $0 \le F_X(x) \le 1$
    \item Non-decreasing function
    \item $\lim_{x \to \infty} F_X(x) = 1$
\end{itemize}

\subsection*{Example}
If $X$ takes values $1,2,3$:
\[
F_X(2) = P(X \le 2) = P(X=1)+P(X=2)
\]

\section*{(6) Probability Density / Mass Function (PDF / PMF)}

\subsection*{Definition}
For a discrete RV, the \textbf{Probability Mass Function (PMF)} is:
\[
p_X(x) = P(X = x)
\]

\textbf{Reasoning:}  
The PMF gives the exact probability at each point and fully characterizes a discrete RV.

\subsection*{Properties}
\begin{itemize}
    \item $p_X(x) \ge 0$
    \item $\sum_x p_X(x) = 1$
\end{itemize}

\subsection*{Example}
For Binomial$(n,p)$:
\[
p_X(k) = \binom{n}{k} p^k (1-p)^{n-k}
\]

This PMF can be used to compute expectation, variance, and probabilities.

\end{document}
